\section{El Agent es una nueva especie}

La sección anterior trató la medición de la inteligencia; esta sección pasa al Agent. Guangmi (广密) llama al Agent «nueva especie» porque ya no es solo un sistema que produce lenguaje, sino un sistema capaz de percibir el entorno, mantener el estado de la tarea, invocar herramientas, ejecutar acciones, leer la retroalimentación y seguir avanzando hacia el objetivo. Cuanto más cerca se esté de la AGI, más podría parecerse la evolución del Agent a un salto de fase que a una mejora lineal.

\begin{figure}[H]
\centering
\includegraphics[width=0.92\textwidth]{agent-three-capabilities.png}
\caption{Las tres capacidades decisivas del Agent: Long context reasoning, Tool use e Instruction following sostienen en conjunto a la nueva especie. Diagrama conceptual propio, elaborado a partir de la conversación de 00:48:03--00:55:49.}
\end{figure}

\begin{knowledgebox}{Lectura de la figura: las tres capacidades son indispensables}
Long Context permite al Agent mantener el estado de las tareas largas, Tool Use le permite modificar el mundo exterior e Instruction Following evita que se aparte del objetivo del usuario. Memory, Feedback y Autonomy son un nivel superior: el Agent no solo ejecuta la orden actual, sino que también acumula experiencia y explora por iniciativa propia.
\end{knowledgebox}

\subsection{Long context reasoning: el estado de las tareas largas}

Esta subsección examina primero el razonamiento con contexto largo. Una tarea de Agent rara vez es una sola pregunta y respuesta: abarca varios pasos, varios archivos, varias herramientas y varias rondas de recuperación de errores. El modelo debe recordar el objetivo, las restricciones, las rutas ya intentadas, las causas de los fallos, los productos actuales y el paso siguiente. El contexto largo es solo el requisito mínimo; la verdadera dificultad está en seleccionar, comprimir y actualizar el contexto.

\begin{warningbox}{El contexto largo no es una memoria universal}
Meter todo en la ventana trae costo, ruido y dilución de la atención. Un buen Agent necesita un estado de tarea estructurado, un sistema de archivos externo, recuperación de información, resúmenes y memoria a largo plazo, no acumular tokens sin límite.
\end{warningbox}

\subsection{Tool use: del lenguaje a la acción}

Esta subsección examina el uso de herramientas. La invocación de herramientas permite al modelo buscar, navegar por páginas web, leer y escribir archivos, ejecutar código, llamar a API y operar interfaces de usuario. Lleva al modelo de «generar sugerencias» a «ejecutar tareas». Pero el uso de herramientas también plantea problemas de seguridad, permisos, recuperación de errores y capacidad de auditoría. Cuanto más capaz es el Agent, más necesita límites explícitos y mecanismos de reversión.

\begin{importantbox}{El contrato básico de un producto de Agent}
El usuario da el objetivo; el Agent descompone la tarea e invoca herramientas; el sistema debe ofrecer control de permisos, visibilidad del avance, explicación de errores, toma de control humana y verificación de resultados. Sin esto, cuanto más capaz sea el Agent, mayor será el riesgo.
\end{importantbox}

\subsection{Instruction following: obediencia ante objetivos complejos}

Esta subsección explica la tercera capacidad. El Agent no solo debe saber usar herramientas, sino también ceñirse de forma sostenida al objetivo del usuario durante tareas de largo alcance: qué no puede hacer, qué recursos no puede usar, qué resultado cuenta como terminado, cuándo debe pedir confirmación y cómo resolver instrucciones contradictorias. El instruction following complejo es lo que permite al Agent pasar de una demo de juguete a un sistema en producción.

\noindent\begin{tabularx}{\textwidth}{p{0.20\textwidth}p{0.34\textwidth}X}
\toprule
Capacidad & Modo de fallo & Requisito de producto \\
\midrule
Razonamiento con contexto largo & Olvida el objetivo, repite intentos, usa información obsoleta & Estado de la tarea y memoria externa. \\
Uso de herramientas & Elige la herramienta equivocada, daña archivos, excede sus permisos & Entorno aislado (sandbox), confirmación, auditoría y reversión. \\
Seguimiento de instrucciones & Se desvía por información intermedia, ignora restricciones & Prioridades explícitas y manejo de conflictos. \\
Aprendizaje a partir de la retroalimentación & Repite el mismo error, no transfiere la experiencia & Registro de fallos, revisión retrospectiva y ciclo cerrado de evaluación. \\
\bottomrule
\end{tabularx}

\subsection{Resumen de la sección}

El Agent es una nueva especie no porque tenga un nombre nuevo, sino porque conecta el modelo de lenguaje con el entorno, las herramientas, la memoria y un ciclo cerrado de acción. La verdadera competencia pasará de la «calidad de la respuesta» a la «calidad con que se completa la tarea».
